\section{Metodolog\'ia, dise\~no e implementaci\'on propuesta}

\subsection{Problema que resolver\'a}
El problema central de esta tesis es la falta de mecanismos pr\'acticos y evaluables para gobernar agentes basados en modelos de lenguaje grandes cuando interact\'uan con herramientas externas y producen efectos en tiempo de ejecuci\'on. En la pr\'actica, un agente no s\'olo genera texto: decide qu\'e herramienta invocar, en qu\'e orden, con qu\'e argumentos y cu\'ando detenerse. Esa autonom\'ia introduce riesgos de seguridad, cumplimiento y trazabilidad, por ejemplo aprobaciones no autorizadas, creaci\'on de tickets sin evidencia suficiente, escritura de memoria fuera de alcance o env\'io de notificaciones sensibles sin aprobaci\'on previa.

El problema no se resuelve \'unicamente con alineaci\'on del modelo, filtrado de respuestas o evaluaci\'on est\'atica de prompts. Se requiere un mecanismo externo capaz de inspeccionar propuestas de acci\'on, validar precondiciones y postcondiciones, bloquear acciones inseguras y permitir recuperaci\'on cuando sea posible. Por ello, la tesis propone un marco experimental para responder, con evidencia reproducible, si los contratos de agente y un gobernador de ejecuci\'on pueden reducir efectos inseguros sin destruir la utilidad operativa del agente.

\subsection{Descripci\'on general de la soluci\'on propuesta}
La soluci\'on propuesta es un \textbf{benchmark de gobernanza en tiempo de ejecuci\'on para agentes con uso de herramientas}, basado en tres elementos: un agente LLM, un conjunto controlado de herramientas deterministas y un \textit{Agent Governor} que aplica contratos formales durante la ejecuci\'on. El agente toma decisiones reales de llamada de herramientas; el gobernador observa cada propuesta y la eval\'ua contra reglas declarativas; y el benchmark mide, en distintos modos de enforcement, qu\'e tan bien se previenen efectos inseguros y qu\'e tanto se preserva la finalizaci\'on segura de tareas.

La propuesta se implementa en cuatro modos experimentales: \texttt{no\_contract}, \texttt{advisory}, \texttt{guarded} y \texttt{strict}. \texttt{no\_contract} funciona como l\'inea base sin intervenci\'on; \texttt{advisory} detecta y registra violaciones pero no bloquea; \texttt{guarded} bloquea acciones inseguras y permite replanteamiento; y \texttt{strict} bloquea y aborta. Esta estructura permite observar no s\'olo si se detectan violaciones, sino tambi\'en el compromiso entre seguridad y utilidad. La arquitectura general del flujo puede resumirse en la Figura~\ref{fig:figure}: el agente propone una acci\'on sobre una herramienta, el Governor eval\'ua esa propuesta contra el contrato activo, decide si la acci\'on se permite, se bloquea o se aborta, y s\'olo despu\'es de esa decisi\'on puede ocurrir el efecto sobre el estado o sobre la herramienta. A continuaci\'on, el runtime registra la traza, actualiza el ledger y entrega la evidencia al evaluador experimental.

\subsection{Descripci\'on detallada por fases y m\'odulos}

\subsubsection{Fase 1: modelado del problema y definici\'on contractual}
La primera fase consiste en traducir riesgos operativos a restricciones observables de ejecuci\'on. En esta tesis, dichas restricciones se expresan como contratos de agente que contienen huellas de configuraci\'on, herramientas declaradas, reglas de pre-ejecuci\'on, reglas de tiempo de ejecuci\'on y reglas de post-ejecuci\'on. Esta fase integra el conocimiento del marco te\'orico sobre gobernanza, control de acceso, validaci\'on por pol\'iticas y auditor\'ia de agentes, con el fin de convertir principios abstractos en reglas verificables.

El resultado de esta fase es un conjunto de contratos YAML que formalizan propiedades como autorizaci\'on, evidencia requerida, alcance permitido de memoria, aprobaci\'on previa y consistencia del ledger de ejecuci\'on. Esta formalizaci\'on constituye una aportaci\'on relevante porque traslada la discusi\'on te\'orica sobre seguridad de agentes a objetos ejecutables y medibles.

\subsubsection{Fase 2: dise\~no del benchmark experimental}
La segunda fase dise\~na el benchmark como un entorno controlado pero suficientemente expresivo para inducir decisiones riesgosas. Para ello se definen escenarios nominales y adversariales, un conjunto peque\~no de herramientas deterministas y un perfil de modelo reproducible. Los escenarios no se eval\'uan contra una secuencia fija de pasos, sino contra \textit{outcomes} aceptables y prohibidos. Esta decisi\'on metodol\'ogica es importante porque evita sobreajustar la evaluaci\'on a una sola trayectoria y permite medir recuperaci\'on segura, rechazo seguro o cumplimiento seguro con distintas rutas de ejecuci\'on.

En esta fase tambi\'en se determinan las expectativas runtime de cada escenario, las propiedades terminales y los casos que deben excluirse de la medici\'on de F1 cuando no representan una oportunidad runtime v\'alida. Con ello, el benchmark mantiene coherencia metodol\'ogica entre dise\~no del escenario, oportunidad de violaci\'on y definici\'on de m\'etricas.

\subsubsection{Fase 3: implementaci\'on del runtime gobernado}
La tercera fase implementa el ciclo de ejecuci\'on del agente. Este runtime construye el contexto del modelo, expone s\'olo las herramientas permitidas por el escenario, recibe propuestas de acci\'on, consulta al gobernador, ejecuta herramientas permitidas, almacena trazas, y sintetiza resultados finales. La implementaci\'on incluye adaptadores de modelo, una capa de herramientas simuladas, c\'alculo de propiedades finales, soporte para \textit{replan}, y mecanismos de \textit{feedback} del gobernador para orientar la recuperaci\'on despu\'es de un bloqueo.

Una aportaci\'on concreta de esta tesis en esta fase es el soporte de recuperaci\'on gobernada: no s\'olo se registra una violaci\'on, sino que el sistema puede impedir la acci\'on insegura y empujar al agente hacia una secuencia v\'alida. El caso de notificaciones sensibles tras aprobaci\'on muestra que la gobernanza eficaz requiere tanto bloqueo como se\~nal de recuperaci\'on.

\subsubsection{Fase 4: instrumentaci\'on, evaluaci\'on y diagn\'ostico}
La cuarta fase implementa la observabilidad experimental. Cada corrida produce \texttt{trace.jsonl}, \texttt{run\_ledger.json} y \texttt{summary.json}. A partir de esos artefactos se calculan m\'etricas de prevenci\'on, finalizaci\'on segura, oportunidad de acci\'on insegura, detecci\'on runtime y sobrecarga. Entre las m\'etricas principales se encuentran \texttt{unsafe\_side\_effect\_rate}, \texttt{governance\_effectiveness}, \texttt{successful\_safe\_completion\_rate}, \texttt{precision}, \texttt{recall} y \texttt{f1}.

Adicionalmente, se incorpora una etapa de diagn\'ostico de F1 para distinguir errores de alineaci\'on del \textit{oracle}, escenarios con baja presi\'on adversarial, y casos estructuralmente excluidos de la medici\'on runtime. Esta parte es importante porque evita conclusiones enga\~nosas y fortalece la validez interna del experimento.

\subsubsection{Fase 5: ejecuci\'on comparativa y an\'alisis}
La quinta fase ejecuta la matriz experimental completa sobre todos los escenarios y modos. El objetivo es comparar la l\'inea base sin contrato contra variantes con observaci\'on, bloqueo con recuperaci\'on y bloqueo con aborto. La hip\'otesis principal es que \texttt{guarded} ofrecer\'a una mejor relaci\'on entre seguridad y utilidad que \texttt{strict}, al prevenir efectos inseguros sin penalizar de forma excesiva la capacidad de completar tareas.

Esta fase culmina con el an\'alisis comparativo de resultados, donde se contrasta cu\'ando la detecci\'on es suficiente, cu\'ando el bloqueo mejora la seguridad, y cu\'ando el aborto reduce de manera innecesaria la utilidad del agente. As\'i, el benchmark no s\'olo mide si una regla se activa, sino si la propuesta completa resulta operativamente conveniente.

\subsection{Aportaci\'on de la tesis y sustento te\'orico}
La aportaci\'on principal de esta tesis es un marco reproducible para estudiar gobernanza externa de agentes LLM con uso de herramientas en tiempo de ejecuci\'on. Esta aportaci\'on se sustenta en tres niveles. Primero, retoma del marco te\'orico la necesidad de mecanismos verificables de control, trazabilidad y cumplimiento para sistemas aut\'onomos. Segundo, toma del estado del arte la observaci\'on de que los agentes modernos ya no pueden evaluarse s\'olo como generadores de texto, sino como sistemas que ejecutan acciones y producen efectos. Tercero, propone una metodolog\'ia experimental concreta para vincular ambos planos mediante contratos ejecutables, instrumentaci\'on reproducible y m\'etricas orientadas a seguridad y utilidad.

En t\'erminos metodol\'ogicos, la tesis no se limita a proponer una arquitectura conceptual; tambi\'en demuestra factibilidad de implementaci\'on mediante un \textit{software artifact} ejecutable, escenarios controlados, integraci\'on con backends LLM reales y evidencia cuantitativa. Esto vuelve sustentable la propuesta, porque la misma puede ser replicada, auditada y extendida a otros modelos, proveedores o entornos con herramientas reales.

\subsection{Calidad argumentativa y congruencia con la investigaci\'on previa}
La propuesta mantiene congruencia con la investigaci\'on desarrollada hasta ahora porque parte del supuesto, respaldado por el estado del arte, de que los modelos con capacidad agentiva requieren gobernanza ex\'ogena cuando interact\'uan con sistemas externos. Tambi\'en es congruente con el marco te\'orico sobre control basado en pol\'iticas, seguridad por restricciones y validaci\'on de transiciones de estado. En lugar de afirmar que la alineaci\'on interna del modelo es suficiente, la tesis argumenta que se necesita una capa complementaria de enforcement y evidencia.

Finalmente, la propuesta es factible porque ya cuenta con implementaci\'on operativa, corpus de escenarios, perfiles de modelo, contratos, diagn\'ostico y evaluaci\'on. En ese sentido, la tesis avanza de una discusi\'on conceptual hacia una demostraci\'on experimental medible. La referencia a la Figura~\ref{fig:figure} permite integrar esta explicaci\'on textual con la representaci\'on visual de la arquitectura, reforzando la coherencia entre dise\~no, implementaci\'on y evaluaci\'on.
